\section{Question 1}
\begin{frame}
    \frametitle{問題設定}
    まず，問題で与えられている集合と関数を整理する．

    \begin{block}{与えられているもの}
        \begin{align*}
            &A = \{2^k \mid k \in \mathbb{N}_0,\ k \le 7\}\\
            &B = \{n \mid n \in \mathbb{N}_0,\ n \le 255\}\\
            &f : \mathcal{P}(A) \to B\\
            &f(X) = \sum_{a \in X} a
        \end{align*}
    \end{block}

    \begin{itemize}
        \item \(\mathbb{N}_0=\mathbb{N}\cup\{0\}\)
        \item \(\mathcal{P}(A)\) は \(A\) のべき集合
        \item \(f(X)\) は部分集合 \(X\) の要素をすべて足した値
    \end{itemize}
\end{frame}

\begin{frame}
    \frametitle{(1) 集合 \(A\) の全要素}
    定義をもう一度見ると，
    \[
        A=\{2^k \mid k \in \mathbb{N}_0,\ k \le 7\}
    \]
    である．ここで \(\mathbb{N}_0=\{0,1,2,3,\dots\}\) なので，
    \[
        k=0,1,2,3,4,5,6,7
    \]
    を代入すればよい．したがって，

    \pause

    \begin{align*}
        2^0 &= 1, &
        2^1 &= 2, &
        2^2 &= 4, &
        2^3 &= 8,  &
        2^4 &= 16, &
        2^5 &= 32, &
        2^6 &= 64, &
        2^7 &= 128
    \end{align*}

    \pause

    \begin{alertblock}{解答}
        \[
            A=\{1,2,4,8,16,32,64,128\}
        \]
    \end{alertblock}
\end{frame}


\begin{frame}
    \frametitle{(2) べき集合 \(\mathcal{P}(A)\) の要素数}
    (1) より，集合 \(A\) の要素は 8 個である．

    \begin{block}{べき集合の要素数}
        要素数が \(n\) の集合 \(S\) について，
        \[
            |\mathcal{P}(S)|=2^n
        \]
        が成り立つ．
    \end{block}

    \pause

    今回は \(|A|=8\) なので，
    \[
        |\mathcal{P}(A)|=2^8=256
    \]
    となる．

    \begin{alertblock}{解答}
        \[
            |\mathcal{P}(A)|=256
        \]
    \end{alertblock}
\end{frame}

\begin{frame}
    \frametitle{(3) \(f(\{1,4\})\) の値}
    関数 \(f\) は，部分集合の要素をすべて足す写像である．

    \begin{block}{関数の定義}
        \[
            f(X)=\sum_{a\in X} a
        \]
    \end{block}

    \pause

    \(X=\{1,4\}\) を代入すると，
    \[
        f(\{1,4\})=1+4=5
    \]
    である．

    \begin{alertblock}{解答}
        \[
            f(\{1,4\})=5
        \]
    \end{alertblock}
\end{frame}


\begin{frame}
    \frametitle{(4) \(f(X)=10\) となる部分集合 \(X\)}
    \(f(X)=10\) とは，部分集合 \(X\) の要素の和が 10 になるという意味である．

    \begin{block}{使える要素}
        \[
            A=\{1,2,4,8,16,32,64,128\}
        \]
    \end{block}

    \pause

    16 以上の要素は 10 より大きいため使えない．
    そこで \(1,2,4,8\) の中から和が 10 になる組合せを探す．

    \pause

    \[
        10=8+2
    \]

    \begin{alertblock}{解答}
        \[
            X=\{2,8\}
        \]
    \end{alertblock}
\end{frame}
